% Auto-generated by scripts/latex_synthesis.py -- do not hand-edit; regenerate from the
% source synthesis/<slug>.md instead, since edits here will be overwritten on the next run.
%
% Thesis-chapter style: report class (numeric \cite by default, no natbib/biblatex needed),
% standard margins, no title page/TOC -- this is meant to be read as a single chapter, not a
% standalone thesis. See scripts/latex_synthesis.py's CITATION REPLACEMENT STRATEGY comment
% for how the "Author et al. (Year)" text was turned into \cite{} calls below.
\documentclass[12pt]{report}
\usepackage[margin=1in]{geometry}
\usepackage[utf8]{inputenc}
\usepackage[T1]{fontenc}

\begin{document}

\chapter{how do these papers handle attention mechanisms}

% Cross-paper synthesis generated by scripts/synthesize.py.

Attention mechanisms emerge across these papers as a solution to the representational bottleneck of compressing variable-length sequences into fixed-size vectors, with later work generalizing attention from an auxiliary alignment tool into the sole computational primitive of the architecture. \cite{1409.3215} exemplifies the problem attention was designed to solve: their sequence-to-sequence LSTM encodes an entire source sentence into one fixed-dimensional vector, a design that works only because of an ad hoc trick—reversing the source sentence—to shorten effective dependency distances, and the paper explicitly notes this fixed-vector bottleneck as a limitation motivating subsequent work \cite{1409.3215}. \cite{1409.0473} directly addresses this bottleneck by introducing a soft-alignment "attention" mechanism that lets the decoder compute a weighted context vector over all encoder hidden states at each output step, rather than relying on a single fixed-length summary; this frees the encoder from having to compress all sentence information into one vector and yields the largest gains on long sentences, precisely the cases where the fixed-vector approach of \cite{1409.3215} struggles. \cite{1706.03762} then takes this alignment idea to its logical extreme, dispensing with recurrence and convolution entirely and building a model, the Transformer, based solely on scaled dot-product and multi-head self-attention; where \cite{1409.0473} uses attention as an add-on to an RNN encoder-decoder, \cite{1706.03762} uses self-attention as the mechanism by which every position in a sequence directly attends to every other position, achieving constant maximum path length between any two positions and much greater parallelism. \cite{1810.04805} builds directly on this architecture, adopting the Transformer's bidirectional self-attention as BERT's core encoder and contrasting it explicitly with the constrained, left-to-right self-attention used in GPT-style unidirectional Transformers, showing that allowing every token to attend freely in both directions (via the masked-language-model objective) yields substantially better representations for downstream token- and sentence-level tasks \cite{1810.04805}.

\cite{10-1371-journal-pone-0130140} and \cite{1512.03385} do not address attention mechanisms in this sense: the former proposes layer-wise relevance propagation for explaining classifier decisions in feedforward and Bag-of-Words models, and the latter introduces residual connections for training very deep image classification networks, so both are omitted from the comparison above.

\bibliographystyle{plain}
\bibliography{how-do-these-papers-handle-attention-mechanisms}

\end{document}
